\chapter{内容要求}
\section{论文题目}
题目应该简短、明确、有概括性。通过题目，能大致了解论文内容、专业特点和学科范畴。但字数要适当，一般不宜超过20字。必要时可加副标题。

\section{摘要与关键词}
\subsection{论文摘要}
摘要应概括反映出毕业设计（论文）的内容、方法、成果和结论。摘要中一般不宜使用公式、图表，不标注引用文献编号。中文摘要以300左右字为宜、外文摘要以250个实词左右为宜。

\subsection{关键词}
关键词是供检索用的主题词条，应采用能覆盖设计（论文）主要内容的通用技术词条（参照相应的技术术语标准），尽量从《汉语主题词表》中选用，未被词表收录的新学科、新技术中的重要术语和地区、人物、文献等名称，也可作为关键词标注。关键词一般为3～5个，按词条的外延层次排列（外延大的排在前面）。关键词应以与正文不同的字体字号编排在摘要下方。多个关键词之间用分号分隔。中英文关键词应一一对应。

\section{目录}
目录按章、节、条三级标题编写，要求标题层次清晰。目录中的标题要与正文中标题一致。目录中应包括绪论、报告（论文）主体、结论、致谢、参考文献、附录等。

\section{正文}
正文是毕业设计(论文)的主体和核心部分，一般应包括绪论、报告（论文）主体及结论等部分。

\subsection{绪论}
绪论一般作为第一章，是毕业论文(设计)主体的开端。绪论应包括：毕业论文（设计）的背景及目的；国内外研究状况和相关领域中已有的成果；设计和研究方法；设计过程及研究内容等。绪论一般不少于1.5千字。

\subsection{主体}
主体是毕业论文(设计)的主要部分，应该结构合理、层次清楚、重点突出、文字简练、通顺。主体的内容应包括以下各方面：
\begin{enumerate}
	\item 毕业论文(设计)总体方案设计与选择的论证。	
	\item 毕业论文(设计)各部分（包括硬件与软件）的设计计算。
	\item 试验方案设计的可行性、有效性以及试验数据的处理及分析。
	\item 对本研究内容及成果进行较全面、客观的理论阐述，应着重指出本研究内容中的创新、改进与实际应用之处。理论分析中，应将他人研究成果单独书写，并注明出处，不得将其与本人的理论分析混淆在一起。对于将其他领域的理论、结果引用到本研究领域者，应说明该理论的出处，并论述引用的可行性与有效性。
	\item 自然科学的论文应推理正确，结论清晰，无科学性错误。
	\item 管理和人文学科的论文应包括对所研究问题的论述及系统分析、比较研究，模型或方案设计，案例论证或实证分析，模型运行的结果分析或建议、改进措施等。
\end{enumerate}

\subsection{结论}
结论是毕业论文(设计)的总结，是整篇设计报告（论文）的归宿。要求精炼、准确地阐述自己的创造性工作或新的见解及其意义和作用，还可进一步提出需要讨论的问题和建议。

\section{参考文献}
按正文中出现的顺序列出直接引用的主要参考文献。

毕业论文(设计)的撰写应本着严谨求实的科学态度，凡有引用他人成果之处，均应按论文中所出现的先后次序列于参考文献中。并且只列出正文中以标注形式引用或参考的有关著作和论文。一篇论著在论文中多处引用时，在参考文献中只能出现一次，序号以第一次出现的位置为准。

\section{致谢}
致谢中主要感谢导师和对论文工作有直接贡献及帮助的人士和单位。

\section{附录}
对于一些不宜放入正文中、但作为毕业论文(设计)又是不可缺少的部分，或有重要参考价值的内容，可编入毕业论文(设计)的附录中。例如，过长的公式推导、重复性的数据、图表、程序全文及其说明等。